\begin{table}[htbp]
\centering
\caption{Data-determined breaks in the U.S. fiscal response}\label{tab:breaks}
\small
\begin{tabularx}{\textwidth}{>{\raggedright\arraybackslash}p{0.3\textwidth}>{\centering\arraybackslash}X>{\centering\arraybackslash}X}
\toprule
Surplus, regressor & Breaks: one / two / three & Number of breaks (sequential / BIC) \\
\midrule
Actual (NIPA, 1971–2024), debt & 1992 / 1992, 2013 / 1979, 1992, 2013 & 1 / 3 \\
Actual, interest & 2009 / 1992, 2009 / 1983, 1995, 2012 & 0 / 3 \\
Cyclically adjusted (CBO, 1967–2024), debt & 1993 / 1979, 1993 / 1979, 1996, 2012 & 1 / 3 \\
Cyclically adjusted, interest & 2009 / 1996, 2009 / 1979, 1996, 2009 & 3 / 3 \\
\bottomrule
\end{tabularx}
\par\smallskip{\footnotesize\raggedright \emph{Notes.} The test of no break against the best number of breaks (UDmax) has a bootstrap p-value of 0.00 in every row.\par}
\end{table}
